% ===========================================================================
% figs/comparators.tex -- fig:comparators, reworked for thesis_v2.
% FIGURE BODY ONLY. No preamble, no document environment. \input it from
% inside a figure/fullwidth float in Sections/04b-representation.tex.
%
% Four instruments, four deliberately incommensurable value-to-position maps,
% one comparator each. The premise and all four scales are v1's and unchanged.
% What changed is labelling only; no mark moved, no number changed.
%
%   (1) ROW B plotted seven marks and named one ("all seven layers"), leaving
%       six unidentifiable. Every mark is now named by the layer it is, with a
%       fanned leader for the three that crowd at 3.6/3.8/4.2, and the stub
%       carries "one mark per layer" so a bare accent numeral above the axis
%       cannot be misread as a second set of tick labels (ticks are quietink
%       and below the axis; layer numbers are accent and above it).
%   (2) ROW A's "slab projected out" sat to the right of its own point and
%       nearer the shuffled-label rule than to the mark it names. It now has a
%       leader back to the 0.121 dot.
%   (3) ROW C's single-cell marker was a hollow ring sitting on the dotted zero
%       rule (-0.0012 against 0.000: 0.10cm apart on this scale). It is now
%       solid, the zero rule is interrupted across its height so nothing runs
%       into it, and its label starts 0.30cm clear with a leader through the
%       interruption -- necessary because the zero rule lies between mark and
%       label and would otherwise claim the label.
%   (4) Bar labels in ROWS C and D moved from a 0.13cm to a 0.30cm gap, so
%       "optimal-transport barycentric" no longer reads as a continuation of
%       its own interval bar.
%   General rule applied throughout: a label more than ~0.3cm from its mark
%   gets a short leader.
%   Row heights grew (rows now 0.70/0.77/0.98 apart rather than colliding) to
%   pay for the new labels. Nothing is rescaled: 1cm here is 1cm on the page.
%
% PROVENANCE of every number drawn is in comparators.json, sibling to this
% file. All four rows were checked against RESULTS_cluster and against the
% thesis text; they agree. Sources:
%   row A, B  RESULTS_cluster/workspace_probe.json
%   row C     RESULTS_cluster/probe_arm_residual.json,
%             probe_rep_residual_s43.json, probe_rep_residual_s45.json,
%             probe_arm_single_cell.json
%   row D     RESULTS_cluster/stratify_sameplate_final.json,
%             stratify_ot_T2.json, stratify_consensus.json
% ===========================================================================
\begin{tikzpicture}[x=1cm, y=1cm, font=\small]
  % Four deliberately different value-to-position maps, one per rule.
  \def\xa#1{{(#1)*10.5}}
  \def\xb#1{{(#1)/20*10.5}}
  \def\xc#1{{((#1)+0.02)/0.13*10.5}}
  \def\xd#1{{((#1)+0.03)/0.08*10.5}}
  % Same maps shifted by a label gap; the shift MUST sit inside the braces --
  % \xa{v}+0.05 is not a pgfmath expression and does not compile.
  % 0.30cm, not v1's 0.13cm: at 0.13cm a label abutted the end of its own
  % interval bar and read as part of it.
  \def\xcR#1{{((#1)+0.02)/0.13*10.5+0.30}}
  \def\xdR#1{{((#1)+0.03)/0.08*10.5+0.30}}
  % leader/label shims: +0.05 and +0.06 clear a 1.4pt marker, -0.10 and +0.10
  % keep row A's two abutting labels off their own dots, +0.35 puts row C's
  % single-cell label 0.30cm clear of the marker's edge.
  \def\xaN#1{{(#1)*10.5+0.05}}
  \def\xaL#1{{(#1)*10.5-0.10}}
  \def\xaP#1{{(#1)*10.5+0.10}}
  \def\xcN#1{{((#1)+0.02)/0.13*10.5+0.06}}
  \def\xcM#1{{((#1)+0.02)/0.13*10.5+0.35}}

  % row dividers
  \foreach \y in {-0.70,-2.76,-5.82}{
    \draw[rulegrey!55, line width=0.3pt] (-4.45,\y) -- (11.20,\y);}

  % ===== ROW A -- is it there? decode accuracy, 0 to 1 ====================
  \draw[ink, line width=0.4pt] (0,0) -- (10.5,0);
  \foreach \v in {0,0.25,0.5,0.75,1}{
    \draw[rulegrey, line width=0.3pt] (\xa{\v},0) -- (\xa{\v},-0.10);
    \node[anchor=north, font=\scriptsize, text=quietink] at (\xa{\v},-0.14) {\v};}
  \draw[ink, line width=0.5pt, densely dotted] (\xa{0.083},-0.14) -- (\xa{0.083},0.52);
  \node[anchor=south west, font=\scriptsize\itshape, text=ink]
    at (\xa{0.083},0.52) {shuffled labels};
  % 0.121 -- the constant post-ablation floor. Its label is pushed clear of the
  % shuffled-label rule 0.40cm to its left and leadered back to the mark.
  \fill[quietink] (\xa{0.121},0) circle (1.4pt);
  \draw[quietink, line width=0.3pt] (\xaN{0.121},0.05) -- (1.57,0.20);
  \node[anchor=west, font=\scriptsize, text=quietink]
    at (1.62,0.22) {slab projected out};
  \fill[quietink] (\xa{0.821},0) circle (1.4pt);
  \node[anchor=south east, font=\scriptsize, text=quietink]
    at (\xaL{0.821},0.08) {raw probe, peak};
  \fill[accent] (\xa{0.883},0) circle (1.4pt);
  \node[anchor=south west, font=\scriptsize, text=accent]
    at (\xaP{0.883},0.08) {context removed};
  \node[anchor=east, align=right, font=\scriptsize, text=ink] at (-0.35,0)
    {\textbf{Is it there?}\\ \textcolor{quietink}{decode accuracy}};

  % ===== ROW B -- is it used? ablation cost in multiples of the null =======
  \draw[ink, line width=0.4pt] (0,-2.05) -- (10.5,-2.05);
  \foreach \v in {0,5,10,15,20}{
    \draw[rulegrey, line width=0.3pt] (\xb{\v},-2.05) -- (\xb{\v},-2.15);
    \node[anchor=north, font=\scriptsize, text=quietink] at (\xb{\v},-2.19) {\v};}
  \draw[ink, line width=0.5pt, densely dotted] (\xb{1},-2.19) -- (\xb{1},-1.35);
  \node[anchor=south west, font=\scriptsize\itshape, text=ink]
    at (\xb{1},-1.35) {matched-dimension random subspace};
  % Seven marks, one per layer, and every mark is named. Four are named
  % directly. The other three are layers 16, 12 and 9, reading 3.6, 3.8 and
  % 4.2: 1.1mm and 2.2mm apart on this axis. Three separate leaders into that
  % converge and stop being readable, and fanning them wide enough to separate
  % would claim a resolution the marks do not have, so the cluster is braced
  % and named as an ordered set instead. Reading order is the axis order.
  \foreach \v in {3.6,3.8,4.2,5.7,12.7,14.1,18.6}{
    \fill[accent] (\xb{\v},-2.05) circle (1.4pt);}
  \draw[accent, line width=0.3pt, decorate,
        decoration={brace, amplitude=2.5pt}] (\xb{3.6},-1.93) -- (\xb{4.2},-1.93);
  \node[anchor=south, font=\scriptsize, text=accent, inner sep=1pt]
    at (\xb{3.9},-1.83) {16, 12, 9};
  \foreach \x/\lay in {5.7/6, 12.7/2, 14.1/8, 18.6/4}{
    \node[anchor=south, font=\scriptsize, text=accent, inner sep=1pt]
      at (\xb{\x},-1.83) {\lay};}
  \node[anchor=east, align=right, font=\scriptsize, text=ink] at (-0.35,-2.05)
    {\textbf{Is it used?}\\
     \textcolor{quietink}{ablation cost, $\times$ the null,}\\
     \textcolor{quietink}{marks named by layer}};

  % ===== ROW C -- is it WHICH drug? fraction of the model's own gap ========
  \draw[ink, line width=0.4pt] (0,-5.05) -- (10.5,-5.05);
  % The label is carried separately from the coordinate so the negative tick
  % sets a real minus and not the ASCII hyphen \foreach would print. Same idiom
  % as gate.tex and power.tex; positives stay unmathed so their digits match
  % rows A and B, which are set in text.
  \foreach \v/\t in {-0.02/$-0.02$, 0/0, 0.02/0.02, 0.04/0.04, 0.06/0.06,
                     0.08/0.08, 0.10/0.10}{
    \draw[rulegrey, line width=0.3pt] (\xc{\v},-5.05) -- (\xc{\v},-5.15);
    \node[anchor=north, font=\scriptsize, text=quietink] at (\xc{\v},-5.19) {\t};}
  % The zero rule is drawn in two pieces. The single-cell median sits at
  % -0.0012, one tenth of a centimetre from zero on this scale, so the rule is
  % interrupted across the marker's height rather than run through it.
  \draw[ink, line width=0.5pt, densely dotted] (\xc{0},-5.15) -- (\xc{0},-4.94);
  \draw[ink, line width=0.5pt, densely dotted] (\xc{0},-4.66) -- (\xc{0},-3.34);
  \draw[ink, line width=0.5pt, densely dotted] (\xc{0.10},-5.15) -- (\xc{0.10},-3.34);
  \node[anchor=south east, font=\scriptsize\itshape, text=ink]
    at (\xc{0.10},-3.34) {pre-registered margin};
  \draw[accent, line width=1.0pt] (\xc{0.0063},-3.60) -- (\xc{0.0335},-3.60);
  \fill[accent] (\xc{0.0197},-3.60) circle (1.4pt);
  \node[anchor=west, font=\scriptsize, text=accent] at (\xcR{0.0335},-3.60)
    {residual arm, anchor seed};
  \draw[quietink, line width=1.0pt] (\xc{0.0060},-4.00) -- (\xc{0.0173},-4.00);
  \fill[quietink] (\xc{0.0118},-4.00) circle (1.4pt);
  \node[anchor=west, font=\scriptsize, text=quietink] at (\xcR{0.0173},-4.00)
    {second seed};
  \draw[quietink, line width=1.0pt] (\xc{0.0029},-4.40) -- (\xc{0.0143},-4.40);
  \fill[quietink] (\xc{0.0085},-4.40) circle (1.4pt);
  \node[anchor=west, font=\scriptsize, text=quietink] at (\xcR{0.0143},-4.40)
    {third seed};
  % Solid, like every other point mark; the absent bar is what says it carries
  % no interval. Leader crosses the interruption in the zero rule.
  \fill[quietink] (\xc{-0.0012},-4.80) circle (1.4pt);
  \draw[quietink, line width=0.3pt] (\xcN{-0.0012},-4.80) -- (\xcR{-0.0012},-4.80);
  \node[anchor=west, font=\scriptsize, text=quietink] at (\xcM{-0.0012},-4.80)
    {single-cell arm, median over its depths};
  \node[anchor=east, align=right, font=\scriptsize, text=ink] at (-0.35,-4.20)
    {\textbf{Is it \emph{which} drug?}\\
     \textcolor{quietink}{fraction of the model's}\\
     \textcolor{quietink}{own preference gap}};

  % ===== ROW D -- does a better target help? scramble gap in NIR ==========
  \draw[ink, line width=0.4pt] (0,-7.40) -- (10.5,-7.40);
  \foreach \v/\t in {-0.02/$-0.02$, 0/0, 0.02/0.02, 0.04/0.04}{
    \draw[rulegrey, line width=0.3pt] (\xd{\v},-7.40) -- (\xd{\v},-7.50);
    \node[anchor=north, font=\scriptsize, text=quietink] at (\xd{\v},-7.54) {\t};}
  \draw[ink, line width=0.5pt, densely dotted] (\xd{0},-7.50) -- (\xd{0},-6.12);
  \draw[accent, line width=1.0pt] (\xd{-0.016},-6.35) -- (\xd{0.042},-6.35);
  \fill[accent] (\xd{0.014},-6.35) circle (1.4pt);
  \node[anchor=west, font=\scriptsize, text=accent] at (\xdR{0.042},-6.35)
    {single cell};
  \draw[accent, line width=1.0pt] (\xd{-0.018},-6.75) -- (\xd{0.016},-6.75);
  \fill[accent] (\xd{-0.001},-6.75) circle (1.4pt);
  \node[anchor=west, font=\scriptsize, text=accent] at (\xdR{0.016},-6.75)
    {optimal-transport barycentric};
  \draw[quietink, line width=1.0pt] (\xd{-0.022},-7.15) -- (\xd{0.003},-7.15);
  \fill[quietink] (\xd{-0.010},-7.15) circle (1.4pt);
  \node[anchor=west, font=\scriptsize, text=quietink] at (\xdR{0.003},-7.15)
    {consensus};
  \node[anchor=east, align=right, font=\scriptsize, text=ink] at (-0.35,-6.75)
    {\textbf{Does a better target help?}\\
     \textcolor{quietink}{scramble gap, $\NIR$,}\\
     \textcolor{quietink}{identifiable stratum}};

\end{tikzpicture}
